\documentclass{BeamerAAV}
\title{Chapitre 4}
\subtitle{Notions pour la statistique multivariée}
\institute{STT-2100 Régression}
\author{Audrey-Anne Vallée}
%\date{Juin 2026}

\usepackage{adjustbox, tikz, hyperref, enumitem}
\usepackage[dvipsnames]{xcolor}



%Raccourcis mathématiques

\newcommand{\esp}{\mathbb{E}}
\newcommand{\var}{\mathbb{V}\text{ar}}
\newcommand{\cov}{\mathbb{C}\text{ov}}
\newcommand{\corr}{\mathbb{C}\text{orr}}
\newcommand{\prob}{\mathbb{P}}
\newcommand{\N}{\mathcal{N}}
\newcommand{\iid}{\overset{iid}{\sim}}


\def\I{\mathbb{I}}
\def\R{\mathbb{R}}
\def\P{{\mathbb P}}


\newcommand{\bm}{\boldmath}
\newcommand{\um}{\unboldmath}
\newcommand{\botheta}{\mbox{\bm$\theta$\um}}
\newcommand{\boTheta}{\mbox{\bm$\Theta$\um}}
\newcommand{\boY}{\mbox{\bm$Y$\um}}
\newcommand{\bobeta}{\mbox{\bm$\beta$}}
\newcommand{\boeps}{\mbox{\bm$\varepsilon$\um}}
\newcommand{\boldx}{\mbox{\bm$x$\um}}
\newcommand{\boldu}{\mbox{\bm$u$\um}}
\newcommand{\boldA}{\mbox{\bm$A$\um}}
\newcommand{\boR}{\mbox{\bm$R$\um}}
\newcommand{\boun}{\mbox{\bm$1$\um}}
\newcommand{\boz}{\mbox{\bm$0$\um}}
\newcommand{\bolambda}{\mbox{\bm$\lambda$\um}}
\newcommand{\bov}{\mbox{\bm$v$\um}}
\newcommand{\boalpha}{\mbox{\bm$\alpha$\um}}
\newcommand{\bo}[1]{\mbox{\bm$#1$\um}}
\newcommand{\boX}{\bo{X}}
\newcommand{\boV}{\bo{V}}
\newcommand{\hbobeta}{\hat{\bobeta}}
\newcommand{\hboeps}{\hat{\boeps}}
\newcommand{\hY}{\hat{\boY}}
\newcommand{\boH}{\bo{H}}
\newcommand{\boI}{\bo{\I}}
\newcommand{\eproof}{\hfill$\sqcap\!\!\!\!\sqcup$}
\newcommand{\bou}{\mbox{\bm$u$\um}}


\begin{document}
\AtBeginSection{\frame{\sectionpage}}

\begin{frame}[plain]
	\titlepage
\end{frame}


\section{Notions d'algèbre linéaire}


\begin{frame}{Définitions et opérations sur les matrices}
\begin{bte}{Matrice}
Une \textbf{matrice} réelle $m\times n$ est un tableau ordonné de nombres réels arrangés en $m$ lignes et $n$ colonnes. 
Une matrice $\boldA$ et ses éléments $a_{ij}$, $i = 1, \dots, m$ et $j=1, \dots, n$, peuvent s'écrire
$$
\boldA=\left[ \begin{array}{ccccc}
a_{11} & \cdots & a_{1j} & \cdots & a_{1n} \\
\vdots &         & \vdots &           & \vdots \\
a_{i1} &  \cdots       &  a_{ij} & \cdots & a_{in} \\
\vdots &         & \vdots &           & \vdots \\
a_{m1} &  \cdots & a_{mj}& \cdots & a_{mn}
\end{array}
\right].
$$
\end{bte}    
\end{frame}

\begin{frame}{Définitions et opérations sur les matrices}
Soit $\boldA$ une matrice de dimension $m\times n$ quelconque.
    \begin{bte}{Transposition}
      La \textbf{transposée} de $\boldA$, notée $\boldA^\top$ ou $\boldA'$, est de dimension $n\times m$ et vérifie $$\alpha_{ji}=a_{ij},$$ $i=1,\ldots,m; \; j=1,\ldots,n.$
    \end{bte}
\begin{exe}{Transposition}
    \[
\boldA=
\begin{bmatrix}
1&2&3\\4&5&6
\end{bmatrix}
\quad\Longrightarrow\quad
\boldA^\top=
\begin{bmatrix}
1&4\\2&5\\3&6
\end{bmatrix}.
\]
\end{exe}
\end{frame}

\begin{frame}{Définitions et opérations sur les matrices}
 \begin{bte}{Symétrie}
      Soit $\boldA$ une matrice de dimension $n\times n$, elle est \textbf{symétrique} si 
    $$ a_{ij} = a_{ji}$$
    pour $i, j = 1, \dots, n$. On a alors que $\boldA = \boldA^\top.$
    \end{bte}
\begin{exe}{Symétrie}
    \[
\boldA=
\begin{bmatrix}
1&2\\2&5
\end{bmatrix}
\quad\Longrightarrow\quad
\boldA^\top=
\begin{bmatrix}
1&2\\2&5
\end{bmatrix}.
\]
\end{exe}
\end{frame}


   
\begin{frame}{Définitions et opérations sur les matrices}
Soit $\boldA$ une matrice carrée de dimension $n\times n$.
    \begin{bte}{Inversion}
      La \textbf{matrice inverse} de $\boldA$, notée $\boldA^{-1}$ est la matrice qui satisfait
    $$ \boldA \boldA^{-1} = \boldA^{-1}\boldA = \bold{I}_n,$$
    où $\bold{I}_n$ est la matrice identité de dimension $n$.
    \end{bte}
    \begin{bte}{Idempotence}
      La matrice $\boldA$ est \textbf{idempotente} si la multiplication par elle-même ne la modifie pas,  
     $\boldA \boldA = \boldA.$
    \end{bte}
\end{frame}

\begin{frame}{Définitions et opérations sur les matrices}
Soient deux matrices $\boldA$ ($m\times n$) et $\bold B$ ($m\times n$), et un scalaire $k\in\R$.
\begin{bte}{Somme matricielle}
    $\bold{C}=\boldA+\bold B$ est la matrice $C$ ($m\times n$) dont les éléments sont
$$c_{ij}=a_{ij}+b_{ij}\qquad\;\;i=1,\ldots,m; \; j=1,\ldots,n.$$
\end{bte}
\begin{bte}{Produit scalaire}
   $\bold{C}=k\boldA$, est telle que
$$c_{ij}=k a_{ij}\qquad i=1,\ldots,m; \; j=1,\ldots,n.$$
\end{bte}
\end{frame}

\begin{frame}{Définitions et opérations sur les matrices}
Soient deux matrices $\boldA$ ($m\times r$) et $\bold B$ ($r\times n$).
\begin{bte}{Produit matriciel}
    $\bold{C}=\bold A\bold B$ correspond à la matrice $\bold C$ ($m\times n$) telle que
$$c_{ij}=\sum_{k=1}^r a_{ik}b_{kj}, \qquad i=1,\ldots,m; \; j=1,\ldots,n.$$
\end{bte}
Notons que le produit matriciel n'est défini que si le nombre de colonnes de $\boldA$ est égal au nombre de lignes de $\bold B$. L'élément $c_{ij}$ correspond au produit scalaire euclidien du i$^\text{ème}$ vecteur ligne de $\boldA$ par le j$^\text{ème}$ vecteur colonne de $\bold B$.
\end{frame}

\begin{frame}{Propriétés du calcul matriciel}
En supposant des matrices compatibles pour les opérations en question, on a que

\begin{enumerate}[label=(\roman*)]
\item Commutativité de la somme : $\boldA + \bold B  =  \bold B + \bold A$
\item Associativité de la somme : $(\bold A+\bold B)+\bold C  =  \bold A+ (\bold B+\bold C)$
\item Associativité du produit:  $(\bold A\bold B)\bold C=\bold A(\bold B\bold C)$
\item Divers types de distributivité :
\begin{align*}
  \bold A(\bold B+\bold C)&=\bold A\bold B+\bold A\bold C \\
(\bold A+\bold B)\bold C&=\bold A\bold C+\bold B\bold C \\
k(\bold A+\bold B)&=k\bold A+k\bold B \\
(k_1+k_2)\bold A&=k_1\bold A+k_2\bold A 
\end{align*}

\item Divers types d'associativité : 
\begin{align*}
    k_1(k_2\bold C)&=k_1k_2\bold C \\
    k(\bold A\bold  B)&=(k\bold A) \bold B=\bold A(k\bold B)
\end{align*}
\end{enumerate}
\end{frame}

\begin{frame}{Propriétés de la transposition et de l'inversion}
\begin{enumerate}[label=(\roman*)]
\item $(\bold A+\bold B)^\top=\bold A^\top+\bold B^\top$
\item $(k\bold A)^\top=k\bold A^\top$
\item $(\bold A\bold B)^\top=\bold B^\top \bold A^\top$
\item $(\bold A^\top)^\top=\bold A$
\item $(\bold A^{-1})^{-1}=\bold A$
\item $(\bold A^\top)^{-1}=\bold (\bold A^{-1})^\top$
\item $(\bold A\bold B)^{-1}=\bold B^{-1} \bold A^{-1}$
\end{enumerate}
\end{frame}

\begin{frame}{Projection orthogonale}
Soit $\bold V$ un sous-espace vectoriel de $\mathbb{R}^n$ et \\
une matrice $\bold X$ dont les colonnes forment une base de $\bold V$. 
\begin{bte}{Matrice de projection orthogonale}
La \textbf{matrice de projection orthogonale} sur l'espace engendré par les colonnes de $\bold X$ est 
$$ \bold P_{\bold X} = \boX \left( \boX^\top \boX\right)^{-1} \boX^\top. $$
$\bold P_{\boX}$  est idempotente et symétrique, $\bold P_{\boX} = \bold P_{\boX}^\top = \bold P_{\boX} \bold P_{\boX}$.
\end{bte}
Soit un vecteur $\bold y \in \mathbb{R}^n$. Sa projection sur $\bold V$ est obtenue le multipliant à la matrice de projection, soit $\bold P_{\boX} \bold y$.
\end{frame}


\begin{frame}{Dériver une fonction d'un vecteur}
\end{frame}
\begin{frame}{Dériver une fonction d'un vecteur}
\end{frame}
\begin{frame}{Dériver une fonction d'un vecteur}
\end{frame}

\begin{frame}{Espérance et variance de variables aléatoires}
\end{frame}

\section{La loi normale multivariée}

\begin{frame}{Loi normale univariée}
\begin{bte}{Normale centrée réduite}
    Soit une variable aléatoire $Z$. Elle est distribuée selon une loi normale centrée réduite, $Z\sim\N(0,1)$, si sa densité s'écrit 
\begin{eqnarray*}
 f_Z(z) = \frac{1}{\sqrt{2\pi}} \exp\left( \frac{-z^2}{2}\right).
\end{eqnarray*}
\end{bte}
\begin{bte}{Normale }
    Une variable aléatoire $X = \sigma Z + \mu$ est distribuée selon une loi normale de moyenne $\mu$ et de variance $\sigma^2$, $X\sim\N(\mu, \sigma^2)$. Sa densité s'écrit 
\begin{eqnarray*}
	f_X(x) =  \frac{1}{\sqrt{2\pi} \sigma} \exp\left\{-\frac{1}{2} \left(\frac{x- \mu}{\sigma}\right)^2\right\}.
\end{eqnarray*}

\end{bte}
\end{frame}

\begin{frame}{Loi normale multivariée}
\begin{bte}{Normale multivariée centrée réduite}
 Soit $\bo{Z}=\left( Z_1, \dots, Z_n \right)^\top$, où $Z_i \overset{iid}{\sim} \N(0,1)$, $i=1, \dots, n$. La densité conjointe des $Z_i$ s'écrit
\begin{eqnarray*}
 f(Z_1, ..., Z_n) = \prod_{i=1}^n f(z_i)
  &=& \left(\frac{1}{\sqrt{2\pi}}\right)^n \exp\left(-\sum_{i=1}^n \frac{z_i^2}{2}\right)\nonumber\\
 &=& \left(\frac{1}{\sqrt{2\pi}}\right)^n \exp\left(-\frac{1}{2} \bo{Z}'\bo{Z}\right)
\end{eqnarray*}
et est la densité d'une loi normale multivariée de dimension $n$, avec $\mathrm{E}(\bo{Z}) = \bo{0}$ et $\mathrm{Var}(\bo{Z}) = \I_n$. On peut écrire
$$\bo{Z}\sim \N_n (\textbf{0}, \boI_n).$$
\end{bte}
\end{frame}


\begin{frame}{Loi normale multivariée}
\begin{bte}{Normale multivariée}
    Soit $\boX=\bo\mu + \bo\Sigma^{(1/2)}\bo{Z}$, où $\bo\mu = (\mu_1,\dots, \mu_n)^\top$ et $\bo\Sigma$ est la matrice de variance-covariance de $\bo X$.
    La densité conjointe des $X_i$ est 
\begin{eqnarray*}
f(X_1, ..., X_n) &=&  \left(\frac{1}{\sqrt{2\pi}}\right)^n \frac{1}{\left|\bo\Sigma\right|^{1/2}} \exp\left\{-\frac{1}{2} (\boX- \bo\mu)' \bo\Sigma^{-1} (\boX- \bo\mu)\right\}.
\end{eqnarray*}
Cette densité est celle d'une loi normale multivariée de dimension $n$, et on peut écrire
\begin{eqnarray*}
\bo{X}\sim \N_n \left(\bo\mu, \bo\Sigma\right) .
\end{eqnarray*}
\end{bte}
\end{frame}

\begin{frame}{Loi normale multivariée}
On a posé $\boX=\bo\mu + \bo\Sigma^{(1/2)}\bo{Z}$. 
On peut déterminer l'espérance et la variance de $\boX$ à partir de celles de $\bo{Z}$.
\vspace{6cm}
\end{frame}

\begin{frame}{Loi normale multivariée}
Il est possible de transformer une variable aléatoire normale. Voici un cas particulier d'intérêt pour le cours.

\begin{bte}{Théorème 1.}
Si 
\begin{itemize}
    \item[-] $\boX \sim \N_n(\bo 0,\boI_n)$
    \item[-]  $\bo A$ est une matrice symétrique
\end{itemize}
alors 
$$\boX'\bo A\boX \sim \chi^2_r$$
si et seulement si $\bo A$ est idempotente de rang r.
\end{bte}
\end{frame}



\end{document}

 
 \begin{bte}{Variable aléatoire}
Une 
 \end{bte}